\chapter{Технологическая часть}
%<писать все про реализацию: какие языки, какие ide и тд.; коды алгоритмов/программы>
%<тестовые данные и волшебные слова <<все тесты пройдены успешно>> >

\section{Реализация алгоритмов}
Для реализации алгоритмов использовался язык программирования C++, так как он поддерживает статическую типизацию и имеет возможность отключения сборщика мусора. Отладка происходила в среде разработки <<Visual Studio 2022>>.\\

\lstinputlisting[firstline=4, lastline=17, style=cpp]{../code/matr.cpp}
Листинг 1 -- Стандартный алгоритм умножения матриц\\
{}

\lstinputlisting[firstline=19, lastline=57, style=cpp]{../code/matr.cpp}
Листинг 2 -- Алгоритм Винограда умножения матриц\\
{}

\lstinputlisting[firstline=59, lastline=91, style=cpp]{../code/matr.cpp}
Листинг 3 -- Алгоритм Винограда умножения матриц с оптимизацией\\
{}

\section{Тестирование}

\begin{table}[!htbp]
	\caption{Модульные тесты для стандартного алгоритма.}
	\centering
	\begin{tabularx}{\textwidth}{|*{2}{>{\centering\arraybackslash}X|}}
		\hline
		\textbf{Вход} & \textbf{Выход}\\\hline
		4 2 2 4 & \text{} \\
		$\begin{pmatrix}
			2&3\\4&5\\3&4\\5&6 
		\end{pmatrix}$
		$\begin{pmatrix}
			3&4&2&5\\6&4&3&2
		\end{pmatrix}$ &
		$\begin{pmatrix}
			24&20&13&16\\42&36&23&30\\33&28&18&23\\51&44&28&37
		\end{pmatrix}$
		\\\hline
	\end{tabularx}
\end{table}

\begin{table}[!htbp]
	\caption{Модульные тесты для алгоритма Винограда.}
	\centering
	\begin{tabularx}{\textwidth}{|*{2}{>{\centering\arraybackslash}X|}}
		\hline
		\textbf{Вход} & \textbf{Выход}\\\hline
		4 2 2 4 & \text{} \\
		$\begin{pmatrix}
			2&3\\4&5\\3&4\\5&6 
		\end{pmatrix}$
		$\begin{pmatrix}
			3&4&2&5\\6&4&3&2
		\end{pmatrix}$ &
		$\begin{pmatrix}
			24&20&13&16\\42&36&23&30\\33&28&18&23\\51&44&28&37
		\end{pmatrix}$ \\\hline
		4 3 3 4 & \text{} \\
		$\begin{pmatrix}
			2&3&3\\4&5&5\\3&4&4\\5&6&6 
		\end{pmatrix}$
		$\begin{pmatrix}
			3&4&2&5\\6&4&3&2\\7&3&2&1
		\end{pmatrix}$ &
		$\begin{pmatrix}
			45&29&19&19\\77&51&33&35\\61&40&26&27\\93&62&40&43
		\end{pmatrix}$
		\\\hline
	\end{tabularx}
\end{table}

\begin{table}[!htbp]
	\caption{Модульные тесты для алгоритма Винограда с оптимизацией.}
	\centering
	\begin{tabularx}{\textwidth}{|*{2}{>{\centering\arraybackslash}X|}}
		\hline
		\textbf{Вход} & \textbf{Выход}\\\hline
		4 2 2 4 & \text{} \\
		$\begin{pmatrix}
			2&3\\4&5\\3&4\\5&6 
		\end{pmatrix}$
		$\begin{pmatrix}
		 	3&4&2&5\\6&4&3&2
		\end{pmatrix}$ &
		$\begin{pmatrix}
			24&20&13&16\\42&36&23&30\\33&28&18&23\\51&44&28&37
		\end{pmatrix}$ \\\hline
		4 3 3 4 & \text{} \\
		$\begin{pmatrix}
			2&3&3\\4&5&5\\3&4&4\\5&6&6 
		\end{pmatrix}$
		$\begin{pmatrix}
			3&4&2&5\\6&4&3&2\\7&3&2&1
		\end{pmatrix}$ &
		$\begin{pmatrix}
			45&29&19&19\\77&51&33&35\\61&40&26&27\\93&62&40&43
		\end{pmatrix}$
		\\\hline
	\end{tabularx}
\end{table}

% чет нечет

\section*{Вывод}
%<что делали и получили в результате>

В данном разделе мной были реализованы алгоритмы умножения матриц и проведено их успешное тестирование.

\clearpage
